\documentclass{situationist-preprint}

\title{On the Word \enquote{Merely}}
\msid{ETH-21-a}
\msfamily{ETH-21}
\msdiscipline{Normative and Applied Ethics}
\mstype{Principal manuscript}
\mskeywords{Kant; formula of humanity; mere means; using persons; consent; reasonable rejection}

\begin{document}
\maketitle

\begin{abstract}
Kant's formula of humanity requires that we treat humanity, in our own person and in that of another, always at the same time as an end and never merely as a means. It is often read as forbidding the use of persons, and so cited against acts in which one person's agency or body serves another's purpose. We argue that we may treat people as means provided we also treat them as ends, and that the whole moral weight of the formula lies in the word \enquote{merely}. We give an account of the qualifier: an agent treats someone merely as a means when she uses him in a way he could reasonably reject, where a ground for rejection is reasonable only if it does not depend on his pursuing a wrongful end. The account explains why ordinary transactions are permissible, why deception and coercion of the innocent are wrong, and why deceiving someone in order to frustrate his wrongdoing need not be. We reply to Parfit's objection that treating someone merely as a means is neither necessary nor sufficient for wrongdoing.
\end{abstract}

\section{The problem}

The second formulation of the categorical imperative reads, in substance: act so that you use humanity, whether in your own person or in that of another, always at the same time as an end, never merely as a means \citep[4:429]{Kant1785}. The formula is cited in applied ethics far more often than any other part of Kant's moral philosophy, and it is often cited as if it said that persons may not be used. Medical research, organ markets, surrogacy and the Footbridge case of the trolley problem have all been condemned on the ground that they \enquote{use} someone.

Read without its qualifier, the formula would forbid most of ordinary life. When I buy bread, I use the baker as a means to obtaining bread; when I ask a stranger for directions, I use her knowledge as a means to finding my way. Kant plainly did not mean to forbid these acts, and he wrote \enquote{merely}. The question is what the qualifier adds. If it adds nothing determinate, the formula reduces to the prohibition of whatever is otherwise wrong, and it does no work \citep{Parfit2011}. If it adds something, we need to say what.

The main accounts in the literature are two. On the \emph{possible consent} account, one treats another merely as a means when one acts on a maxim to which he cannot possibly consent, as in deception and coercion, where the victim is not in a position to dissent \citep{ONeill1985,Korsgaard1986}. On the \emph{attitudinal} account, one treats another merely as a means when one regards him as a mere instrument, so that one would treat him in whatever way best served one's aim \citep{Parfit2011}. We argue that each captures part of the truth and that neither is adequate as it stands, and we propose an account on which the formula forbids exactly the uses that the person used could reasonably reject.

\section{Framework}

\begin{definition}[Use]\label{d:use}
An agent $X$ \emph{uses} a person $Y$ in acting on a maxim $M$ if $M$'s end is to be brought about, in part, by means of $Y$'s body, agency, beliefs or abilities.
\end{definition}

\begin{definition}[Reasonable rejection]\label{d:reject}
A person $Y$ \emph{could reasonably reject} a way of being used if he has a ground for objecting to it that does not depend on his pursuing a wrongful end, and that is not outweighed by the grounds others have for objecting to the alternatives.
\end{definition}

\begin{definition}[Merely as a means]\label{d:merely}
$X$ treats $Y$ \emph{merely as a means} in acting on $M$ if $X$ uses $Y$ in acting on $M$ and $Y$ could reasonably reject the way $M$ uses him. $X$ treats $Y$ \emph{as an end} in acting on $M$ if $M$ uses $Y$, if at all, only in ways that $Y$ could not reasonably reject.
\end{definition}

Definition~\ref{d:reject} adapts Scanlon's contractualist test to the formula \citep{Scanlon1998}; a related proposal is made by \citet{Kerstein2013}. The second clause of Definition~\ref{d:merely} makes treating as an end and treating merely as a means contradictories on uses: every use of a person is one or the other. This is the reading on which the formula's two clauses express a single requirement, as Kant's \enquote{at the same time} suggests.

\section{The argument}

\begin{proposition}\label{p:ordinary}
In ordinary transactions to which the parties consent, each uses the other as a means and treats the other as an end.
\end{proposition}

\begin{proof}
The buyer uses the baker's labor, and the baker the buyer's money; so each uses the other. Each has consented to the way he is used, on terms he knows. A person may object to being used in a way to which he has given informed and uncoerced consent only on grounds he has waived, and a waived ground is not a ground for reasonable rejection. So neither could reasonably reject the use, and by Definition~\ref{d:merely} each treats the other as an end.
\end{proof}

\begin{proposition}\label{p:innocent}
Deceiving or coercing a person who is not pursuing a wrongful end, in order to use him for an end he has not adopted, treats him merely as a means, unless the grounds others have for objecting to the alternatives outweigh his.
\end{proposition}

\begin{proof}
A person who is not pursuing a wrongful end has a ground for objecting to deception or coercion that does not depend on any wrongful end: that his agency is directed at ends he has not chosen and in a way that bypasses his judgment. This is the ground the possible consent account identifies \citep{ONeill1985}. By Definition~\ref{d:reject} the ground supports reasonable rejection unless it is outweighed, and by Definition~\ref{d:merely} the use then treats him merely as a means.
\end{proof}

\begin{proposition}\label{p:wrongdoer}
Deceiving a person in order to frustrate his wrongful end, by no more than is needed to frustrate it, uses him without treating him merely as a means.
\end{proposition}

\begin{proof}
The person's only ground for objecting to the deception is that it frustrates his end, and that ground depends on his pursuing a wrongful end; by Definition~\ref{d:reject} it does not support reasonable rejection. Any further imposition, beyond what frustrating his end requires, he could reasonably reject, since his objection to it would not depend on his wrongful end; the condition \enquote{no more than is needed} excludes such impositions.
\end{proof}

\begin{theorem}\label{t:main}
It is consistent with the formula of humanity to treat a person as a means, provided he is at the same time treated as an end. What the formula forbids is fixed entirely by the qualifier \enquote{merely}.
\end{theorem}

\begin{proof}
By Definition~\ref{d:merely}, a use of a person is either a use he could reasonably reject or one he could not. The formula forbids the first and, by its second clause, requires the second. Proposition~\ref{p:ordinary} shows that the uses it permits include the ordinary cooperative uses of persons, so that using is not as such forbidden; Propositions~\ref{p:innocent} and~\ref{p:wrongdoer} show that whether a deceptive use is forbidden depends on whether the person used could reasonably reject it, which is the condition expressed by \enquote{merely}. Deleting the qualifier would forbid every use, including those of Proposition~\ref{p:ordinary}; so the content of the formula, beyond the prohibition of every use, lies wholly in the qualifier.
\end{proof}

Proposition~\ref{p:wrongdoer} is the result on which ETH-04 relies: the householder who lies to the murderer at the door uses him, but she does not treat him merely as a means, since his only objection to the lie depends on his intention to kill.

\section{Objections and replies}

\begin{objection}[Parfit's gangster]\label{o:parfit}
\citet{Parfit2011} describes an egoist who buys coffee and pays for it, but who would harm the seller if that served his ends. The egoist treats the seller merely as a means, since he regards her as a mere instrument, yet he does nothing wrong in buying the coffee. Treating someone merely as a means is therefore not sufficient for wrongdoing.
\end{objection}

\begin{reply}
The objection relies on the attitudinal account, on which treating merely as a means is a matter of the agent's dispositions across counterfactual situations. On Definition~\ref{d:merely}, it is a matter of the maxim on which the agent acts. The egoist's maxim is to buy coffee at the posted price, and that use of the seller is one she could not reasonably reject. He does not, in this act, treat her merely as a means, though he is a person who would. Parfit's case shows a defect of character that is not a wrong in the act, and the formula, as a principle for the assessment of acts, should not count it as one.
\end{reply}

\begin{objection}[Necessity]\label{o:necessity}
Some acts are wrong without using anyone, as when one lets a person drown to avoid getting wet. The formula, so read, cannot explain their wrongness.
\end{objection}

\begin{reply}
We agree, and Theorem~\ref{t:main} does not claim otherwise. The formula's prohibition concerns uses of persons. The duty to aid, which Kant treats as an imperfect duty following from the same formula's positive requirement to further others' ends, explains the drowning case; ETH-18 argues that the duty to rescue holds up to moderate cost. The formula is not a complete theory of the right, and nothing in the word \enquote{merely} suggests that it is.
\end{reply}

\begin{objection}[Circularity]\label{o:circ}
Definition~\ref{d:reject} uses the notion of a wrongful end. If wrongness is what the formula is meant to explain, the account is circular.
\end{objection}

\begin{reply}
The circle is not vicious. The definition refers to the wrongness of the end pursued by the person used, not to the wrongness of the act being assessed, and the former can be settled independently, often by the formula itself applied to the other person's act. The structure is recursive, as is usual for principles that assign forfeiture of claims to those who violate them.
\end{reply}

\begin{objection}[Trolley cases]\label{o:trolley}
In the Footbridge case the man is used as a means, and the wrongness of pushing him is the formula's paradigm application. If using is permitted when the person is also treated as an end, the account must explain why he is treated merely as a means, and it is not clear that it can.
\end{objection}

\begin{reply}
The man on the footbridge has a ground for objecting to being pushed that does not depend on any wrongful end, namely that he is subjected to a lethal threat he was not under, and the account classifies the use as one he could reasonably reject. But we do not hold that the formula is what distinguishes Footbridge from Switch. ETH-03 explains the distinction by the doctrine of the redirected threat and shows that the means principle, which forbids every use without consent, gives the wrong verdict in the Loop case. Our reading of the formula is consistent with that result: the one on the loop is used, but the formula forbids use only when it is mere use, and whether a use is mere use depends on the claims of the person used, which ETH-03 identifies.
\end{reply}

\section{What the result does not establish}

The result does not establish which uses persons could reasonably reject in contested cases such as markets in organs or commercial surrogacy; it shows only that the question of their permissibility is not settled by the observation that someone is used. Nor does it bear on whether the constraint expressed by the formula is better stated in deontological or in consequentialist form. ETH-02 argues that the choice between these forms is notational, and the constraint, whatever its content, can be stated in either. The present argument concerns its content.

\section{Conclusion}

We may treat people as means. The formula of humanity forbids treating them merely as means, and it forbids nothing else. On the account defended here, a person is treated merely as a means when he is used in a way he could reasonably reject, and he can reasonably reject only on grounds that do not depend on his own wrongful ends. The moral weight of the formula lies in the word \enquote{merely}.

\begin{thebibliography}{99}\small
\bibitem[Kant(1785)]{Kant1785} Kant, I. (1785). \emph{Groundwork of the Metaphysics of Morals}. In \emph{Practical Philosophy}, trans. M.~J. Gregor, Cambridge University Press, 1996.
\bibitem[Kerstein(2013)]{Kerstein2013} Kerstein, S.~J. (2013). \emph{How to Treat Persons}. Oxford University Press.
\bibitem[Korsgaard(1986)]{Korsgaard1986} Korsgaard, C.~M. (1986). The right to lie: Kant on dealing with evil. \emph{Philosophy \& Public Affairs}, 15(4), 325--349.
\bibitem[O'Neill(1985)]{ONeill1985} O'Neill, O. (1985). Between consenting adults. \emph{Philosophy \& Public Affairs}, 14(3), 252--277.
\bibitem[Parfit(2011)]{Parfit2011} Parfit, D. (2011). \emph{On What Matters}, Volume One. Oxford University Press.
\bibitem[Scanlon(1998)]{Scanlon1998} Scanlon, T.~M. (1998). \emph{What We Owe to Each Other}. Harvard University Press.
\end{thebibliography}

\end{document}
